\chapter{METHODOLOGY}

\label{ch:methodology}

\section{Research Methodology}

The work follows the Design Science Research Methodology \cite{hevner2004}, in which an artifact is
built to address a problem class, demonstrated against the problem, and evaluated with evidence
appropriate to the artifact. For a planning tool the appropriate evidence is not a satisfaction
survey but a demonstration that the artifact enforces the properties it claims and degrades
explicitly where it cannot.

\subsection{Iterations}

The thesis was executed as a sequence of small, independently reviewable increments rather than as
a single build. Each increment had to satisfy four conditions before it counted as delivered: it
added one coherent capability; its new tests failed against the code they replace; it passed a
single continuous-integration gate; and it recorded what it deliberately did not do. The
increments grouped into five phases.

\begin{enumerate}
  \item \textbf{Foundation.} Accounts, sessions, and the global course catalogue.
  \item \textbf{Program context.} The program-aware data model, an imported reference curriculum,
        and the rewrite of every progress, grade and rating read to resolve per program.
  \item \textbf{Numeric grades and difficulty.} Immutable attempts, a credit-weighted average, and
        completion-gated ratings with a shrinkage estimate.
  \item \textbf{Planning.} Workload analysis, semester recommendations, the saved-plan browser, and
        the atomic reference-selection editor with drafts, blockers and ordering.
  \item \textbf{Administration.} Resource settings, cohort demand, capacity diagnostics, two
        allocation models, an explicit job executor, and immutable captures with history.
\end{enumerate}

\subsection{Evaluation strategy}

Evaluation is confined to three questions the artifact can actually answer. Does it enforce its
stated invariants? The answer is a test suite that exercises the invariants against a real database,
including real concurrent writers and real transaction rollback. Can the result be reproduced? The
answer is a pinned runtime, a committed lockfile, a disposable database, and one continuous-
integration job that runs the whole gate on every change. Does it remain usable by people who do not
read the code? The answer is a browser accessibility audit and a measured profiling run, both of
which are reported with their limits rather than as headline numbers.

What the evaluation deliberately does not attempt is an outcome study. Recommending courses is a
decision-support problem with a long feedback loop and many confounders; measuring whether the
recommendation engine improves graduation outcomes would require a population, a control condition
and a multi-year horizon, none of which were available.

\section{Requirements}

The problem statement in Chapter~\ref{ch:introduction} yields functional requirements, each paired
with the property that makes it testable.

\begin{enumerate}
  \item \textbf{Program context.} A course may belong to several programs; placement and
        prerequisites resolve inside one program. Testable property: no reader falls back to a
        global catalogue or to a union of programs when a context is present.
  \item \textbf{Account scope.} Every mutation carries the owner and program context the client
        observed. Testable property: a mismatched scope is refused with a conflict, and the stored
        data is unchanged.
  \item \textbf{Authoritative completion.} Completion validates prerequisites, and uncompletion
        cascades transitively. Testable property: after a cascade, no remaining completion has an
        unmet prerequisite.
  \item \textbf{Immutable grade history.} A retake adds an attempt; it does not overwrite one.
        Testable property: an identical retry returns the original attempt; a changed payload for
        the same key is refused.
  \item \textbf{Honest estimates.} A difficulty estimate from few votes is pulled toward a prior
        and always displayed with its vote count.
  \item \textbf{Atomic saved selections.} A saved semester is validated, canonicalised, revision
        checked and committed with its receipt in one transaction.
  \item \textbf{Frozen simulations.} An allocation model that changes is a new version, not an
        edit, and a stored capture replays under the version that produced it.
  \item \textbf{Explicit absence.} Missing data produces an explicit unknown state, never a zero,
        a default or a fabricated total.
\end{enumerate}

\section{Architecture}

The system is an npm-workspace monorepo with three packages. The {\em shared} package holds Zod
schemas and inferred TypeScript types and is compiled before either application so that both sides
import the same contract. The {\em server} package is an Express application over Prisma and
PostgreSQL. The {\em client} package is a React single-page application built with Vite.

\begin{table}[h]
\centering
\small
\begin{tabular}{p{0.10\textwidth}p{0.44\textwidth}p{0.32\textwidth}}
\toprule
\textbf{Package} & \textbf{Responsibility} & \textbf{Verification} \\
\midrule
shared & Zod contracts, DTOs, frozen simulation schemas & compiled first, typechecked \\
server & REST API, transactions, services & Jest on real PostgreSQL \\
client & routes, panels, guards, stores & Vitest in jsdom \\
root   & scripts, CI, startup supervision & one Ubuntu job \\
\bottomrule
\end{tabular}
\caption{Package responsibilities and their verification surface}
\end{table}

Three architectural rules are load-bearing rather than stylistic.

\textbf{The server is authoritative.} The client may optimistically pre-compute a cascade so the
interface responds immediately, but the stored result is whatever the server's transaction decided,
and the client reconciles to the returned snapshot.

\textbf{Every write is scoped.} A mutation carries the expected owner and program context. The
server compares them against the values it reads inside the same transaction that performs the
write, so a client holding a stale view cannot commit against a changed context.

\textbf{Absence is a value.} Unknown credits, unknown seats, unknown grade eligibility and
unresolved requirements have their own representation. A missing resource configuration is not zero
capacity; a course with no votes is not a course of average difficulty; a program with no verified
credit total is not a program of zero credits.

\section{Data Model}

The schema has twenty-three models and six enumerations. Six of them carry the program context and
four carry the simulation.

\begin{table}[h]
\centering
\small
\begin{tabular}{p{0.44\textwidth}p{0.44\textwidth}}
\toprule
\textbf{Model} & \textbf{Role} \\
\midrule
{\tt Course} & global identity; never duplicated per program \\
{\tt Curriculum} & program metadata; credit total may be unknown \\
{\tt CurriculumCourse} & membership; unique per program and course \\
{\tt CurriculumPlacement} & one source occurrence with year, semester and group \\
{\tt CurriculumPrerequisite} & a prerequisite edge, naming a member or an outside course \\
{\tt CurriculumCohort} & the admission year a program structure governs \\
{\tt CurriculumRequirement} & an uncoded requirement such as a free elective \\
{\tt User} & account, role, optional program assignment \\
{\tt StudentRecord} & completion or plan per course, with elective claim \\
{\tt GradeAttempt} & immutable attempt with an owner-scoped key \\
{\tt CourseRating} & one vote per account and course \\
{\tt StudyPlan} & saved plan metadata \\
{\tt PlannedSemester} & stored selection, revision and reference metadata \\
{\tt ReferenceSelectionReceipt} & immutable receipt for an atomic save \\
{\tt SchoolResource} & scenario settings with a compare-and-set revision \\
{\tt SimulationAllocationRun} & immutable one-course aggregate capture \\
{\tt SimulationSemesterRun} & immutable semester aggregate capture \\
{\tt SimulationAllocationJob} & durable one-course request \\
{\tt SimulationSemesterAllocationJob} & durable semester request \\
{\tt SimulationSemesterParticipant} & private per-student outcome \\
\bottomrule
\end{tabular}
\caption{The models that carry program context, evidence and simulation state}
\end{table}

Membership and placement are separate for a reason. A course appears once in a program but may be
listed in several semesters and in several elective groups: the reference curriculum imports
repeated placements as placement rows distinguished by source order, while membership stays unique.
Prerequisite edges reference the (program, course) pair rather than the course alone, which is what
makes one program's edges independent of another's.

\section{Program Data Processing}

The reference curriculum was obtained from the program's published curriculum page and imported by
a staged pipeline rather than by a direct write. A read-only verifier first loads the source and
proves that nothing in it is lost: course identities must be unique within the source, credits must
equal the sum of their lecture and laboratory units, repeated placements must be preserved, a blank
course code is accepted only for an explicit free-elective requirement, and elective selection
counts must agree within a group. A compatibility preflight then inspects the target database, and
only then does an idempotent {\tt --apply} backfill create the program, its memberships, its
placements, its requirements and its prerequisite edges, recording the source label and order.

The import deliberately does not compute a degree total. The legacy reference's modelled path
totals disagree with the credit total stated in later published versions of the same program, and
no authoritative reconciliation was available. The stored total is therefore null, every reader
treats it as unknown, and the interface reports an unavailable percentage rather than a number
derived from an assumption.

\section{Verification Strategy}

Server tests run against a real PostgreSQL database rather than a mock, because the properties
worth testing here are transactional: cascade behaviour, lock ordering, rollback after a partial
write, and an exact retry that must not duplicate an immutable row. Suites run sequentially against
a shared instance so that fixtures do not fight over predicate locks, while individual race tests
still issue genuinely concurrent requests. Client tests run in a DOM emulator and cover state
machines, stale replies, guards and rendering; the properties that need a real engine -- colours,
focus movement, zoom limits -- are checked in a browser.

The gate is one Ubuntu job: pinned Node and npm, restore from a committed lockfile, build shared,
build both applications, typecheck, lint with zero warnings, check formatting, apply migrations and
seed a disposable service, then run startup, server and client suites. A change that cannot pass
that job is not merged. Where a run failed for a reason that could not be established, the failure
was recorded as an open finding rather than absorbed by relaxing an assertion.
